\pagestyle{empty}
% ======[COVER]======
\begin{center}
%\begin{small}
\textbf{\textsc{a thesis submitted for the degree of}}

\textbf{\textsc{Doctor of Philosophy}}
%\end{small}

\vspace*{\fill}

\begin{Huge}
{\color{oxblue}\textbf{Widening the applicability}}
\end{Huge}

\vspace{2mm}

\begin{Huge}
{\color{oxblue}\textbf{of permutation inference}}
\end{Huge}

\vspace*{\fill}

\textbf{Anderson M.\ Winkler}

\textbf{St.\ Edmund Hall}

\vspace*{\fill}

\includegraphics[scale=1.6]{figures/ox_logo_special_gold_pos.eps}

\vspace*{\fill}

\textbf{\textsc{fmrib} Analysis Group}

\textbf{Nuffield Department of Clinical Neurosciences}

\textbf{University of Oxford}

\vspace*{\fill}

\textbf{Trinity 2016}
\end{center}

% ======[ABSTRACT]======
\cleardoublepage
\setstretch{1}
\renewcommand\abstractname{{\Large Abstract}}
\begin{abstract}
This thesis is divided into three main parts. In the \textsc{\textbf{first}}, we discuss that, although permutation tests can provide exact control of false positives under the reasonable assumption of exchangeability, there are common examples in which global exchangeability does not hold, such as in experiments with repeated measurements or tests in which subjects are related to each other. To allow permutation inference in such cases, we propose an extension of the well known concept of exchangeability blocks, allowing these to be nested in a hierarchical, multi-level definition. This definition allows permutations that retain the original joint distribution unaltered, thus preserving exchangeability. The null hypothesis is tested using only a subset of all otherwise possible permutations. We do not need to explicitly model the degree of dependence between observations; rather the use of such permutation scheme leaves any dependence intact. The strategy is compatible with heteroscedasticity and can be used with permutations, sign flippings, or both combined. In the \textsc{\textbf{second}} part, we exploit properties of test statistics to obtain accelerations irrespective of generic software or hardware improvements. We compare six different approaches using synthetic and real data, assessing the methods in terms of their error rates, power, agreement with a reference result, and the risk of taking a different decision regarding the rejection of the null hypotheses (known as the resampling risk). In the \textsc{\textbf{third}} part, we investigate and compare the different methods for assessment of cortical volume and area from magnetic resonance images using surface-based methods. Using data from young adults born with very low birth weight and coetaneous controls, we show that instead of volume, the permutation-based non-parametric combination (NPC) of thickness and area is a more sensitive option for studying joint effects on these two quantities, giving equal weight to variation in both, and allowing a better characterisation of biological processes that can affect brain morphology.
\end{abstract}

% ======[ACKNOWLEDGEMENTS]======
\cleardoublepage
\newpage
%\setstretch{\lspac}
%\vspace*{\fill}
%\begin{center}
%\begin{Large}
%\textbf{Acknowledgements}
%\end{Large}
%\end{center}

\chapter*{Acknowledgements}
\setstretch{\lspac}

\noindent
\`{A} minha fam\'{i}lia.

\vspace{3mm}\noindent
To my friends, for their infinite patience and support.

\vspace{3mm}\noindent
To my doctoral advisors, Prof.\ Stephen M.\ Smith and Prof.\ Thomas E.\ Nichols, endless sources of knowledge and inspiration.

\vspace{3mm}\noindent
I am very much indebted to the National Research Council of Brazil (Conselho Nacional de Desenvolvimento Cient\'{i}fico e Tecnol\'{o}gico --- \textsc{cnp}q), without which this work would have been impossible.

\vspace{3mm}\noindent
A strong, prolific, and enriching collaboration have helped and improved all the chapters, and I am very much thankful for these fruitful interactions. Chapter~\ref{sec:ptree} received helpful feedback from Matthew A.\ Webster, Diego Vidaurre, and Gerg\H{o} Nemes. Likewise, Chapter~\ref{sec:accel} received much welcome input from Gerard R.\ Ridgway and Gwena\"{e}lle Douaud. Chapter~\ref{sec:cortex} would have been impossible without the close collaboration with Lars M.\ Rimol, Jon Skranes, Knut J\o{}rgen Bjuland, Douglas N.\ Greve, Asta K.\ H{\aa}berg, and Mert R.\ Sabuncu.

\vspace{3mm}\noindent
My time in Oxford was made even better thanks to the people in the colleges to which I was at one time or another affiliated: St.\ Cross and St.\ Edmund Hall. Under the risk unfair omissions, some I would like to thank nominally Prof.\ Keith Gull, Dr.\ Richard Willden. Prof.\ Heidi Johansen--Berg, Prof.\ Paul M.\ Matthews, Dr.\ Charlotte Stagg, Sir Mark Jones, Prof.\ Petros Ligoxygakis, and Dr. Katie Warnaby, all of whom, through their remarkable work and friendliness, have made these years extremely enjoyable.

\vspace{3mm}\noindent
At \textsc{fmrib}, the list is long and, again, in the hope of not unfairly omitting anyone, I am much thankful to the members of the Analysis Group listed here in no particular order: Mois{\'e}s Hern{\'a}ndez--Fern{\'a}ndez, Fidel Alfaro--Almagro, Samuel J.\ Harrison, Jonathan Hadida, Giles L.\ Colclough, Zobair Arya, Emmanuel Vallee, Riham Satti, Gwena\"{e}lle Douaud, Eugene P.\ Duff, Gerard R.\ Ridgway, Diego Vidaurre, Matteo Bastiani, Evangelos Roussos, Michiel Cottaar, Eelke Visser, Stamatios Sotiropoulos, Mark Jenkinson, Jesper L.\ Andersson, Saad Jbabdi, Mark Woolrich, Christian F.\ Beckmann, Janine Bijsterbosch, Oiwi Parker Jones, Matthew A.\ Webster, Sean Fitzgibbon, Paul McCarthy, Ludovica Griffanti, Emma C.\ Robinson, Jelena Bo\v{z}ek Mouthuy, David Flitney, Duncan Mortimer, Thomas E.\ Nichols and Stephen M.\ Smith. Additionally, I would like to thank Susan Field and Marilyn Goulding for being tirelessly solicitous since even before the time I arrived, and Marion Greenleaves for all her support.

\vspace{3mm}\noindent
It would have been far more difficult to do this work without the use of free software. I very much thank every person who contributes, directly or indirectly, to the development of these tools.

\vspace{3mm}\noindent
Finally, I would like to thank the Sir Edward Penley Abraham Fund and St.\ Edmund Hall for the honour of being twice the recipient of the ``Cephalosporin Award'', providing the much needed further support during this period, and the Nuffield Department of Clinical Neurosciences and St.\ Edmund Hall for the funding that facilitated attending and presenting work at conferences.
